% Fig. 2 -- MXFP4 fused-engine datapath (exemplar of the four format-specialised
% engines). Arithmetic drawn as circles (x=multiply, +=add), control/scan stages
% as boxes, matching mxdotp_fused_engine.sv / mx_acc_slide / mx_find_lead /
% mx_finalize exactly (widths verified against the RTL ports).
\begin{tikzpicture}[
  font=\small, >=Stealth,
  mul/.style={circle, draw, minimum size=6.5mm, inner sep=0pt, fill=blue!6},
  add/.style={circle, draw, minimum size=6.5mm, inner sep=0pt, fill=orange!10},
  shiftbox/.style={diamond, draw, aspect=2.4, minimum width=15mm, minimum height=8mm,
                   inner sep=0pt, fill=violet!8},
  stg/.style={draw, rounded corners=1pt, align=center, minimum height=9mm,
              minimum width=17mm, fill=blue!7, font=\small},
  opnd/.style={draw, rounded corners=1pt, align=center, minimum height=7mm,
               font=\small, fill=black!4},
  lbl/.style={font=\scriptsize, align=center},
  fl/.style={-{Stealth[length=5pt]}},
  scl/.style={-{Stealth[length=5pt]}, green!40!black, thick},
]
  %=================== operand row ===================
  \node[opnd, minimum width=17mm] (rs1) at (0,5.0)   {\texttt{rs1}$=A$\\16$\times$E2M1};
  \node[opnd, minimum width=17mm] (rs2) at (3.0,5.0) {\texttt{rs2}$=B$\\16$\times$E2M1};
  \node[opnd, minimum width=26mm] (rs3) at (10.0,5.0) {\texttt{rs3}: $S_a,S_b$ (8\,b ea.), old acc (FP32)};

  %=================== 16-lane multiply-accumulate ===================
  \node[mul] (m0) at (0.4,3.6) {$\otimes$};
  \node[mul] (m1) at (1.6,3.2) {$\otimes$};
  \node[lbl] at (2.35,3.05) {$\cdots$};
  \node[mul] (m15) at (3.1,3.6) {$\otimes$};
  \node[lbl, anchor=north] at (m0.south)  {\shortstack{$a_0$\\$b_0$}};
  \node[lbl, anchor=north] at (m1.south)  {\shortstack{$a_1$\\$b_1$}};
  \node[lbl, anchor=north] at (m15.south) {\shortstack{$a_{15}$\\$b_{15}$}};
  \draw[fl] (rs1) -- (1.4,4.4) -- (m1.north);
  \draw[fl] (rs2) -- (1.6,4.4) -- (m1.north);

  \node[isosceles triangle, draw, shape border rotate=-90, isosceles triangle apex angle=60,
        minimum height=16mm, fill=orange!8] (sumtri) at (5.0,3.6) {};
  \node[lbl] at ($(sumtri.center)+(-0.05,0)$) {$\Sigma$};
  \draw[fl] (m0)  -- ($(sumtri.west)+(0,0.55)$);
  \draw[fl] (m1)  -- (sumtri.west);
  \draw[fl] (m15) -- ($(sumtri.west)+(0,-0.55)$);
  \node[lbl] at (1.85,3.9) {9\,b product};

  \node[stg, minimum width=16mm, fill=blue!14] (sop) at (7.1,3.6) {SoP\\13\,b, $\alpha{=}2$\\resident};
  \draw[fl] (sumtri.east) -- (sop);

  %=================== accumulator slide + frame ===================
  \node[shiftbox] (shft) at (10.0,3.6) {$\gg$};
  \node[lbl, above=0.5mm of shft] {\texttt{mx\_acc\_slide}};
  \draw[fl] (rs3.south) |- (shft.north);

  \node[add] (plus1) at (12.6,3.6) {$+$};
  \draw[fl] (sop.east) -- (plus1.west |- sop.east) -- (plus1.west);
  \draw[fl] (shft.east) -- (plus1.west);

  \node[stg, minimum width=18mm, fill=blue!7] (word) at (12.6,1.6) {63\,b word\\$\{38\text{ frame}\mid25\text{ rem}\}$};
  \draw[fl] (plus1.south) -- (word.north);

  \node[stg, minimum width=17mm] (lzc) at (9.5,1.6) {leading-one\\scan (LZC)};
  \draw[fl] (word.west) -- (lzc.east);

  \node[stg, minimum width=17mm, fill=orange!10] (fin) at (6.5,1.6) {normalise\\+ round};
  \draw[fl] (lzc.west) -- node[lbl,above]{lead pos} (fin.east);

  \node[opnd, fill=orange!18, minimum width=14mm] (res) at (3.7,1.6) {FP32\\result};
  \draw[fl] (fin.west) -- (res.east);

  %=================== scale path (bottom, green) ===================
  \node[add, fill=green!10] (splus) at (10.0,-0.3) {$+$};
  \draw[scl] ($(rs3.south)+(-2mm,0)$) |- ($(splus.north)+(-2mm,0)$);
  \draw[scl] (splus.north) -- (shft.south);
  \draw[scl] (splus.west) -| (fin.south);
  \node[lbl, anchor=north, xshift=6mm] at (splus.south) {$S_a{+}S_b{-}254$\\$=$ scale\_exp (16\,b)};

  \node[lbl, text=green!40!black, align=center] at (6.0,-1.35)
       {scale\_exp is hoisted once in the front-end and consumed twice ---\\
        the shift amount at \texttt{mx\_acc\_slide} and the result exponent at
        \texttt{mx\_finalize} --- never touching the 63-bit word.};
\end{tikzpicture}